\documentclass{article}
\usepackage{times}
\usepackage{a4wide}
\usepackage{latexsym}
\usepackage{mathrsfs}
\usepackage{amsmath}
\usepackage{amssymb}
\usepackage{amsthm}
\usepackage{ifthen}
\usepackage{stmaryrd}
\usepackage{algorithm}
\usepackage{ctex}
%\usepackage{algorithmic}
\usepackage{algpseudocode}
\usepackage{graphics}
\usepackage{epsfig}
\usepackage{setspace}
\setstretch{1}
%\usepackage{ramacros}

\addtolength{\textwidth}{20mm}
\addtolength{\oddsidemargin}{-10mm}
\addtolength{\textheight}{18mm}
\addtolength{\topmargin}{-9mm}


\newcounter{exercise}
\setcounter{exercise}{0}
\newcommand{\exercise}{
        \addtocounter{exercise}{1}
        \vspace{0.2in}
        \noindent
        {\bf \theexercise .}
        }

\newcommand{\REMARK}[1]{}


\newcommand{\NEWPART}{\vspace{.1in}
              \noindent}

\newcommand{\<}{
    \langle}

\renewcommand{\>}{
    \rangle}

\newcommand{\ceil}[1]{\left\lceil #1 \right\rceil}

\pagestyle{plain} \pagenumbering{arabic}

\title{{\bf Assignment 4} \\ {\large Deadline: May 5}}

\author{125033910296 李博洋}
\date{}

\begin{document}
\maketitle

%{\large \noindent
%\begin{tabular}{lcl}
%{Jan 17 2007}.
%\end{tabular}
%}

{\large

\begin{exercise}
	Prove the complementary slackness property of linear programs.
\end{exercise}
\\\\
\noindent\textbf{Solution:}
\\
\textbf{原问题与对偶问题}\\
给定原线性规划（最大化）：
\[
\text{(P)}\quad
\begin{cases}
\max\; \boldsymbol{c}^\mathrm{T}\boldsymbol{x} \\[4pt]
\text{s.t.}\; A\boldsymbol{x} \le \boldsymbol{b},\\[4pt]
\boldsymbol{x} \ge \boldsymbol{0}.
\end{cases}
\]
其对偶规划（最小化）：
\[
\text{(D)}\quad
\begin{cases}
\min\; \boldsymbol{b}^\mathrm{T}\boldsymbol{y} \\[4pt]
\text{s.t.}\; A^\mathrm{T}\boldsymbol{y} \ge \boldsymbol{c},\\[4pt]
\boldsymbol{y} \ge \boldsymbol{0}.
\end{cases}
\]

\noindent\textbf{互补松弛性定理}\\
设 $\boldsymbol{x}$ 为原问题可行解，$\boldsymbol{y}$ 为对偶问题可行解。
则 $\boldsymbol{x}$ 为原问题最优解、$\boldsymbol{y}$ 为对偶问题最优解
\iff 满足互补松弛条件：
\[
\boldsymbol{y}^\mathrm{T}\big(\boldsymbol{b}-A\boldsymbol{x}\big) = 0,\quad
\boldsymbol{x}^\mathrm{T}\big(A^\mathrm{T}\boldsymbol{y}-\boldsymbol{c}\big) = 0.
\]

\noindent\textbf{预备：弱对偶定理}\\
对任意原可行解 $\boldsymbol{x}$ 与对偶可行解 $\boldsymbol{y}$，必有：
\[
\boldsymbol{c}^\mathrm{T}\boldsymbol{x}
\le \boldsymbol{y}^\mathrm{T}A\boldsymbol{x}
\le \boldsymbol{b}^\mathrm{T}\boldsymbol{y}.
\]
\begin{proof}
由 $A^\mathrm{T}\boldsymbol{y}\ge \boldsymbol{c},\boldsymbol{x}\ge\boldsymbol 0$，左乘 $\boldsymbol{x}^\mathrm{T}$ 得
$\boldsymbol{c}^\mathrm{T}\boldsymbol{x} \le \boldsymbol{x}^\mathrm{T}A^\mathrm{T}\boldsymbol{y}$；
由 $A\boldsymbol{x}\le \boldsymbol{b},\boldsymbol{y}\ge\boldsymbol 0$，左乘 $\boldsymbol{y}^\mathrm{T}$ 得
$\boldsymbol{y}^\mathrm{T}A\boldsymbol{x} \le \boldsymbol{b}^\mathrm{T}\boldsymbol{y}$。
结合即证。
\end{proof}

\noindent\textbf{互补松弛性定理证明}\\
\textbf{必要性：最优解 $\Rightarrow$ 互补松弛}\\
设 $\boldsymbol{x}^*$ 为原最优解，$\boldsymbol{y}^*$ 为对偶最优解。
由\textbf{强对偶定理}，原对偶最优目标值相等：
\[
\boldsymbol{c}^\mathrm{T}\boldsymbol{x}^*
= \boldsymbol{b}^\mathrm{T}\boldsymbol{y}^*.
\]
结合弱对偶不等式：
\[
\boldsymbol{c}^\mathrm{T}\boldsymbol{x}^*
=\boldsymbol{y}^{*\mathrm{T}}A\boldsymbol{x}^*
=\boldsymbol{b}^\mathrm{T}\boldsymbol{y}^*.
\]

\begin{enumerate}
\item 由 $\boldsymbol{y}^{*\mathrm{T}}A\boldsymbol{x}^* = \boldsymbol{b}^\mathrm{T}\boldsymbol{y}^*$，移项：
\[
\boldsymbol{y}^{*\mathrm{T}}\big(\boldsymbol{b}-A\boldsymbol{x}^*\big)=0.
\]
\item 由 $\boldsymbol{c}^\mathrm{T}\boldsymbol{x}^* = \boldsymbol{y}^{*\mathrm{T}}A\boldsymbol{x}^*$，转置移项：
\[
\boldsymbol{x}^{*\mathrm{T}}\big(A^\mathrm{T}\boldsymbol{y}^*-\boldsymbol{c}\big)=0.
\]
\end{enumerate}
$(\boldsymbol{x}^*,\boldsymbol{y}^*)$ 满足互补松弛，必要性得证。

\noindent\textbf{充分性：可行+互补松弛 $\Rightarrow$ 最优解}\\
已知：
$\boldsymbol{x}$ 原可行，$\boldsymbol{y}$ 对偶可行，且
\[
\boldsymbol{y}^\mathrm{T}(\boldsymbol{b}-A\boldsymbol{x})=0,\quad
\boldsymbol{x}^\mathrm{T}(A^\mathrm{T}\boldsymbol{y}-\boldsymbol{c})=0.
\]

展开条件：
\[
\boldsymbol{y}^\mathrm{T}\boldsymbol{b} = \boldsymbol{y}^\mathrm{T}A\boldsymbol{x},\quad
\boldsymbol{c}^\mathrm{T}\boldsymbol{x} = \boldsymbol{x}^\mathrm{T}A^\mathrm{T}\boldsymbol{y}.
\]
因此：
\[
\boldsymbol{c}^\mathrm{T}\boldsymbol{x}
=\boldsymbol{y}^\mathrm{T}A\boldsymbol{x}
=\boldsymbol{b}^\mathrm{T}\boldsymbol{y}.
\]

原可行解与对偶可行解目标函数值相等，
由线性规划最优判定准则：$\boldsymbol{x}$ 为原最优解，$\boldsymbol{y}$ 为对偶最优解。
充分性得证。


\begin{exercise}
	Transform following problems into linear programming
	
	(1)
		$$\begin{aligned}
		&\text { maximize } & 5x+2y\\
		&\text { subject to } & 0 \le x \le 20 \\
		& & |x - y|\le 5
		\end{aligned}$$
		
	(2)
		$$\begin{aligned}
		&\text { maximize }  & min(x1,x2,x3)\\
		&\text { subject to }  & x1+x2+x3 = 15
		\end{aligned}$$
\end{exercise}
\\\\
\noindent\textbf{Solution:}
\\
(1) 拆开绝对值不等式
$$\begin{aligned}
&\text { maximize } & 5x+2y &\\
&\text { subject to } & 0 \le x \le 20 &\\
& & x - y \le 5 & \rightarrow x \le y + 5 &\\
& & -5 \le x - y & \rightarrow y + 5 - 10 \le x
\end{aligned}$$
令 $z = y + 5$，则有
$$\begin{aligned}
&\text { maximize } & 5x+2z &\\
&\text { subject to } & x \le 20 &\\
& & x - z\le 0 &\\
& & z - x \le 10 &\\
& & x \ge 0, z \ge 0
\end{aligned}$$
\\
(2) 选择$x_1$作为最小值，则有
$$\begin{aligned}
&\text { maximize } & x_1 &\\
&\text { subject to } & x_1 + x_2 + x_3 = 15 &\\
& & x_1 \le x_2 & \rightarrow x_1 - x_2 \le 0 &\\
& & x_1 \le x_3 & \rightarrow x_1 - x_3 \le 0
\end{aligned}$$


\begin{exercise}
	Given a graph $G$, each vertex $v_i$ has a profit $p_i$ and each edge $e_{ij}$ has a cost $c_{ij}$. Define the profit of a cycle is the total profits of all the vertices in the cycle, and the cost of a cycle is the total costs of all the edges in the cycle. We need to find a simple cycle in $G$ which contains a given vertex $v_0$, and maximize the profit of it within the cost bound $B$. Write the linear programming formulation of this problem.
\end{exercise}
\\\\
\noindent\textbf{Solution:}
\\
\textbf{目标函数（最大化圈总收益）：}
\[
\text { maximize } \quad Z = \sum_{i \in V} p_i x_i
\]

\noindent\textbf{约束条件：}
\begin{align*}
x_0 &= 1 \quad &&\text{必须包含指定顶点} \ v_0 \\
\sum_{j:(i,j) \in E} y_{ij} &= 2x_i, \ \forall i \in V \quad &&\text{圈度数约束，每个被选的点的度数为2} \\
\sum_{(i,j) \in E} c_{ij} y_{ij} &\le B \quad &&\text{总成本预算约束} \\
0 \le x_i &\le 1, \ \forall i \in V \quad &&\text{顶点变量范围} \\
0 \le y_{ij} &\le 1, \ \forall (i,j) \in E \quad &&\text{边变量范围}
\end{align*}


\begin{exercise}
Spectrum management is the process of regulating the use of radio frequencies
to promote efficient use and gain a net social benefit. Given a set of broadcast
emitting stations $s_1, . . . , s_n$, a list of frequencies $f_1, . . . , f_m$, where $m \ge n$, and
the set of adjacent stations ${(s_i
, s_j )}$ for some $i, j \in [n]$. The goal is to assign
a frequency to each station so that adjacent stations use different frequencies
and the number of used frequencies is minimized. Formulate this problem as an
integer linear programming problem.
\end{exercise}
\\\\
\noindent\textbf{Solution:}
\\
\noindent\textbf{决策变量}\\
\begin{itemize}
    \item \( x_{ik} \in \{0,1\} \)：发射台 \( s_i \) 分配频率 \( f_k \) 取1，否则取0，\( i=1,\dots,n,\ k=1,\dots,m \)
    \item \( y_k \in \{0,1\} \)：频率 \( f_k \) 被使用取1，否则取0，\( k=1,\dots,m \)
\end{itemize}
\noindent\textbf{目标函数（最小化使用频率数量）}\\
\[
\min \quad \sum_{k=1}^{m} y_k
\]
\noindent\textbf{约束条件}\\
\begin{align*}
    &\sum_{k=1}^{m} x_{ik} = 1, && \forall i = 1,\dots,n \quad \text{（每个发射台分配唯一频率）} \\
    &x_{ik} + x_{jk} \leq 1, && \forall (s_i,s_j) \text{（相邻台）},\ \forall k = 1,\dots,m \quad \text{（相邻台频率不同）} \\
    &x_{ik} \leq y_k, && \forall i=1,\dots,n,\ \forall k=1,\dots,m \quad \text{（频率使用标记约束）} \\
    &x_{ik}, y_k \in \{0,1\}, && \forall i,k \quad \text{（整数约束）}
\end{align*}


\begin{exercise}
You have been chosen to assign the new graduate students to office spaces. Office $i$ has $n_i$ available desks, and graduate student
$g$ is willing to pay you pig dollars for the privilege of being placed in office $i$.

\begin{itemize}
    \item Show how you can use a min-cost flow algorithm to work out the best beneficial (for you) assignment of new students to offices.
    \item Argue that you can achieve this best result without forcing grad students to share
a desk or split their time between multiple rooms.
    \item Write a concise linear program to solve this problem (assuming it returns an
integral answer).
\end{itemize}

\end{exercise}
\\\\
\noindent\textbf{Solution:}
\\
\textbf{1. 使用最小费用流算法解决最佳分配问题}\\
将最大收益分配问题转化为\textbf{最小费用流}问题（收益取反为负成本），构建流网络如下：
\begin{itemize}
    \item \textbf{节点}：源点 $S$、研究生节点 $G_g$、办公室节点 $O_i$、汇点 $T$。
    \item \textbf{有向边（容量，单位费用）}：
    \begin{enumerate}
        \item $S \to G_g$：容量 $1$，费用 $0$（每个研究生仅分配一次）；
        \item $G_g \to O_i$：容量 $1$，费用 $-p_{ig}$（最大化收益 $\leftrightarrow$ 最小化负成本）；
        \item $O_i \to T$：容量 $n_i$，费用 $0$（办公室 $i$ 的工位容量限制）。
    \end{enumerate}
    \item 求解\textbf{最小费用最大流}，总流量为研究生总人数；总费用的相反数即为最大收益。
\end{itemize}
\textbf{2. 证明无需强制学生共享桌位或分时段使用多个房间}\\
依据最小费用流的\textbf{整数流性质}：
所有边的容量均为整数，因此最优流必然是\textbf{整数流}。
\begin{itemize}
    \item 边 $S \to G_g$ 容量为 $1$：每个研究生仅分配至\textbf{一个办公室}，无时间拆分；
    \item 边 $O_i \to T$ 容量为 $n_i$：每个办公室最多容纳 $n_i$ 人，无工位共享。
\end{itemize}
\textbf{3. 线性规划问题解决该问题}\\
定义决策变量：$x_{gi}=1$ 表示研究生 $g$ 分配到办公室 $i$，$x_{gi}=0$ 表示未分配。
\begin{align*}
    \max &\sum_{g} \sum_{i} p_{ig} x_{gi} && \\
    \text{s.t.} &\sum_{i} x_{gi} = 1, &&\forall \text{ 研究生 } g, (\text{每个研究生仅分配至一个办公室})\\
    &\sum_{g} x_{gi} \leq n_i, &&\forall \text{ 办公室 } i, (\text{每个办公室最多容纳 } n_i \text{ 人})\\
	&x_{gi} \in \{0,1\} &&\forall g,i. (\text{整数约束})
\end{align*}




\begin{exercise}
In a \textsc{Facility Location Problem}, there is a set of facilities and a set of cities, and our job is to choose a subset of facilities to open, and to connect every city to some one of the open facilities. There is a nonnegative cost $f_j$ for opening facility $j$, and a nonnegative connection cost $c_{(i,j)}$  for connecting city $i$ to facility $j$. Given these as input, we look for a solution that minimizes the total cost. Formulate this facility location problem as an integer programming problem.
\end{exercise}
\\\\
\noindent\textbf{Solution:}
\\
\noindent\textbf{符号定义}
\begin{itemize}
    \item 设施集合：$J = \{1,2,\dots,m\}$，城市集合：$I = \{1,2,\dots,n\}$
    \item $f_j \geq 0$：开设设施$j$的固定成本
    \item $c_{ij} \geq 0$：城市$i$连接到设施$j$的连接成本
    \item 决策变量：
    \begin{itemize}
        \item $x_j \in \{0,1\}$：$x_j=1$表示开设设施$j$，否则为0
        \item $y_{ij} \in \{0,1\}$：$y_{ij}=1$表示城市$i$连接设施$j$，否则为0
    \end{itemize}
\end{itemize}

\noindent\textbf{数学模型}\\
\noindent\textbf{目标函数（最小化总费用）：}
\[
\min \quad \sum_{j \in J} f_j x_j + \sum_{i \in I} \sum_{j \in J} c_{ij} y_{ij}
\]

\noindent\textbf{约束条件：}
\begin{align*}
    \sum_{j \in J} y_{ij} &= 1, &&\forall i \in I \quad \text{（每个城市唯一连接）} \\
    y_{ij} &\leq x_j, &&\forall i \in I, j \in J \quad \text{（仅连接已开设设施）} \\
    x_j &\in \{0,1\}, &&\forall j \in J \quad \text{（设施0-1变量）} \\
    y_{ij} &\in \{0,1\}, &&\forall i \in I, j \in J \quad \text{（连接0-1变量）}
\end{align*}


\end{document}
